\documentclass[12pt]{article}

% fonts

\usepackage[T1]{fontenc}
\usepackage{newtxtext}
\usepackage{bm}
\usepackage{amsthm} % before newtxmath b/c latex
\usepackage{newtxmath}
% \usepackage[bb=mt,cal=mt]{mathalpha}

\usepackage{amsmath}
\usepackage{amsfonts}
\usepackage{thmtools}
\usepackage{mathtools}

% spacing

\usepackage[parfill]{parskip}
\setlength{\parindent}{0pt}
\setlength{\parskip}{2ex plus 0.2ex minus 0.2ex}

% geometry of the page

\usepackage[
paper  = letterpaper,
left   = 1in,
right  = 1in,
top    = 1in,
bottom = 1in,
]{geometry}

% references

\usepackage{natbib}


\begin{document}

\begin{flushleft}
\textbf{Probabilistic Topic Models (Blei, 2012)} \\
Satvik Saha (ss7228) \\
\today
\end{flushleft}

\vspace{0.1in}

This article surveys probabilistic topic models, focusing on \emph{Latent
Dirichlet Allocation} (LDA), which can be described by the generative model
\begin{align*}
    \beta_k &\sim \text{Dirichlet}(\eta), \\
    \theta_i &\sim \text{Dirichlet}(\alpha), \\
    z_{ij}\mid \theta_i &\sim \text{Categorical}(\theta_i), \\
    w_{ij}\mid z_{ij}, \beta &\sim \text{Categorical}(\beta_{z_{ij}}),
\end{align*}
where $\eta \in \mathbb{R}^V$, $\alpha \in \mathbb{R}^K$.
We only observe a collection of words $\{w_{ij}\}$, with $w_{ij}$ being the
$i$-th word (from a vocabulary of $V$ words) from the $j$-th document.
The goal is to discover a hidden structure of $K$ topics $\{\beta_k\}$, with
each topic $\beta_k$ being a distribution over words.
Any given document (which itself can be viewed as a distribution over words) is
assumed to be a mixture of these topics, with proportions $\theta_i$; this
makes LDA a kind of \emph{mixed membership model} as opposed to a simple
mixture model.
The hyperparameter $\alpha$ controls the sparsity of $\theta_i$; sparser
$\theta_i$ meaning that fewer topics dominate in the corresponding document.

LDA can be viewed as a kind of factorization model: the observed document-word
distributions are factored into document-topic distributions $\theta$ and
topic-word distributions $\beta$.
A non-probabilistic version would involve matrix factorization, such as in
\emph{Nonnegative Matrix Factorization} (NMF).

LDA can also be viewed as a kind of dimensionality reduction technique: the
inferred document-topic distribution $\theta_i \in \Delta^{K - 1}$ is a summary
of the underlying document-word generating distribution (which lives in
$\Delta^{V - 1}$) when $K \ll V$.

A crude way of incorporating metadata (publishing information, location, etc)
$\{x_i\}$ for each document might be to draw $\theta_i \mid \alpha_i \sim
\text{Dirichlet}(\alpha_i)$, with $\alpha_i$ itself regressed on $x_i$ in some
way.
This might help explore how different factors in $x_i$ push documents towards
particular topics.



\nocite{*}
\bibliographystyle{plainnat}

\vfill

\begingroup
\renewcommand{\section}[2]{}%
\textbf{References}
\bibliography{week05}
\endgroup


\end{document}

%%% Local Variables:
%%% mode: latex
%%% TeX-master: t
%%% End:
